\chapter{第一讲\quad 刑法论}
\label{lecture:01}
\noindent{\small\hyperref[ref:lecture:01]{查阅本讲扩展规则}}\par

\section{第一节\quad 刑法的解释}

\paragraph{教材模型\quad 给词语加上条文下标}
设 \(r\) 为条文，\(w\) 为术语，\(D_r(w)\) 是该词在条文语境下的通常外延，\(B_r(w)\) 是在该语境下有理由支持的可能文义边界，\(E_r(w)\) 是待论证的解释外延。三者都以同一类型的事例为元素，并先约定 \(D_r(w)\subseteq B_r(w)\)。
\begin{align*}
 E_r=D_r &\quad \text{平义},\\
 D_r\subsetneq E_r\subseteq B_r &\quad \text{扩大},\\
 E_r\subsetneq D_r &\quad \text{缩小},\\
 x\notin B_r,\ x\text{与典型事例相似}
 &\quad \text{仅凭相似把 }x\text{ 纳入属于类推}.
\end{align*}
集合公式描述已经提出的解释方案，\(B_r\) 的划定仍须论证。两方案只有交叉、没有包含时，应写明增删的部分，单凭「看起来范围更大」无法归类。反义解释另需证明规定具有排他性，例示性条款无法直接用反义封闭。形式上，从 \(P\Rightarrow Q\) 推出 \(\neg P\Rightarrow\neg Q\) 会犯否定前件的错误。增加 \(Q\Rightarrow P\) 或证明列举已穷尽后，该推论才成立。

\paragraph{理由与外延的关系}
文理理由检查语词与句法，体系理由检查条文位置、并列罪名和兜底项，目的理由检查所保护的法益。同一解释方案可以同时获得多种理由。把不同条文的同名谓词分别写成 \(E_{r_1}(w)\)、\(E_{r_2}(w)\)，就可以理解货币犯罪中并列的「伪造」与「变造」和其他条文的术语不必同范围。解释兜底项「其他方法」时，先从列举项抽取共同性质 \(K\)，再检查 \(K(x)\) 与文义边界。以抢劫罪为例，按照教材的功能解释，需检查方法能否压制被害人反抗，仅能便利取得财物的手段还不足以满足该门槛。

当然解释需要补出一条单调性前提。若 \(x\preceq y\) 表示在\emph{同类、同一构成机制}下 \(y\) 更重，而且已论证该规则对严重程度向上封闭，则
\[
 x\preceq y,\quad Q(x),\quad
 \forall a,b\,((a\preceq b\land Q(a))\Rightarrow Q(b))
 \quad\vdash\quad Q(y).
\]
向上封闭一旦成立，也可由 \(\neg Q(y)\) 反推 \(\neg Q(x)\)。性质不同、各自所需法定要素不同的两个罪名，无法仅凭严重程度使用该规则。

\paragraph{推导\quad 串通拍卖与串通投标}
设事实 \(F_1\) 为仅在拍卖中协调竞买价，\(F_2\) 为无招标投标活动。按教材对第223条的文义分析，「投标」以招标投标活动为语境。由 \(F_2\) 与该语境条件推出本例不符合该条的投标要素。\(F_1\) 所表现的竞争失真可以论证规制目的相近，却无法补足投标要素。依第3条，以第223条给本例入罪的推导在该要素处停止。若案件事实上存在招标投标活动，应依新事实重新判断，而不是复用前例结论。

\paragraph{反例与观点题}
以物的财产属性解释抢劫罪，无法因为行为严重就把婴儿纳入「财物」。形式错误出在比较关系的定义域。主观题如给出两种有据可辩的解释，先写共同事实 \(F\)，再分别写
\[
  F+\mathcal R_1\vdash Q,\qquad F+\mathcal R_2\vdash\neg Q,
\]
并点明换掉的是保护目的、文义边界或其他解释前提。教材标注的学说分布属于教材信息，答题还须核对题设指定立场及现行规范。

\section{第二节\quad 刑法的原则}

\paragraph{法条与审查条件}
第3条要求依法定罪处刑，第4条要求平等适用，第5条要求刑罚与罪行及刑事责任相适应。设 \(\mathrm{Basis}(g,t)\) 表示时点 \(t\) 有适用的法定罪刑依据，\(\mathrm{Match}(g,x)\) 表示案件 \(x\) 符合罪名 \(g\) 的必要要件，\(\mathrm{Convict}(g,x)\) 表示依法可作该罪的有罪结论。抽取第3条约束可写为
\[
 \mathrm{Convict}(g,x)
 \ \Longrightarrow\ \mathrm{Basis}(g,t)\land\mathrm{Match}(g,x).
\]
这是必要条件。查到一条法条，尚须做事实归摄、排除事由和证明审查。教材以「成文、事前、严格、确定」概括罪刑法定的审查方向。成文检查罪刑依据与制定权限，事前检查生效时点及第12条例外，严格检验文义边界，确定检验构成与法定刑能否合理预见。行政规范可以参与认定刑法已规定的规范要素，单凭行政规范不能增设犯罪或刑罚。

\paragraph{推导\quad 民事责任与刑事责任可并存}
令 \(\mathrm{Tort}(x)\) 为满足民事侵权要件，\(\mathrm{Injury}(x)\) 为在题设给定条件下达到故意伤害罪的构成门槛。假设甲故意殴打乙，乙轻伤，伤害归因于殴打，具备责任条件，已排除防卫且证明充分。先将故意、伤害结果与因果事实分别代入第234条相应要素，得 \(\mathrm{Injury}(x)\)。再将过错、侵害、损害和因果事实代入民事侵权规则，得 \(\mathrm{Tort}(x)\)。两条证明链在部分事实处重叠，因而可以同时推出 \(\mathrm{Tort}(x)\land\mathrm{Injury}(x)\)。
若把结果换成仅轻微伤，并设未触及其他犯罪或犯罪形态，民事结论仍可成立，原入罪链却缺少结果门槛。于是 \(\mathrm{Tort}\Rightarrow\mathrm{Injury}\) 与 \(\mathrm{Tort}\Rightarrow\neg\mathrm{Injury}\) 都没有普遍有效性。

\paragraph{反例\quad 法定区间与量刑理由}
设适用档内法定主刑范围为 \(I_g\)，量刑事实集为 \(S\)。先从已认定的犯罪情节确定 \(I_g\)，再依 \(S\) 及从轻、从重、减轻等规则确定刑罚。若某案在普通幅度内，提出「被告人身份特殊所以越过法定上限」既无合法升档依据，也无法由第5条推出。第4条要求同样事实适用同样规则，差别待遇须有可指认的法定条件。第5条提供理由约束，无从推出对每项情节赋固定分数的通用量刑函数。

\section{第三节\quad 刑法的效力}

\paragraph{空间效力\quad 第6至11条}
令 \(A\) 表示有犯罪行为发生在我国领域内，\(R\) 表示有犯罪结果发生在我国领域内，\(V\) 表示犯罪发生在我国船舶或航空器内，\(E\) 表示适用法定例外。属地常规入口写作 \(((A\lor R)\land\neg E)\lor V\)，并另核船舶、航空器案件是否涉及特别规则。以 \(C\) 表示中国公民，\(N\) 表示国家工作人员或军人，\(G\) 表示行为为我国刑法规定之罪，\(H\) 表示境外犯罪，\(U,L\) 分别表示所考察的法定最高刑与最低刑。主要入口和法律措辞可分开记：
\begin{align*}
 H\land C\land G
   &\Rightarrow \text{适用本法},\\
 H\land C\land\neg N\land G\land(U\le3\text{年有期徒刑})
   &\Rightarrow \text{可以不予追究},\\
 H\land N\land G
   &\Rightarrow \text{依第7条第2款适用本法},\\
 H\land\mathrm{Foreign}\land\mathrm{AgainstCN}
 \land G\land(L\ge3\text{年有期徒刑})\land\mathrm{LocPun}
   &\Rightarrow \text{依第8条可以适用本法}.
\end{align*}
这里 \(\mathrm{Foreign}\) 为外国人，\(\mathrm{AgainstCN}\) 为对我国国家或公民犯罪，\(\mathrm{LocPun}\) 为依犯罪地法律受处罚。国家工作人员、军人分支不用第7条第1款的轻罪例外。第9条需查我国缔结或参加的条约、义务范围及本法的罪名依据，不能只靠行为的国际性质。第10条所涉外国审判后仍可依法追究，已在外国受刑的可免除或减轻处罚。外交特权和豁免的刑事责任按第11条经外交途径解决。

\paragraph{推导\quad 上限和下限相反}
外国人在境外对中国公民实施行为，按题设依两地法律均受处罚。我国适用刑档为二年至五年有期徒刑。于是 \(L=2<3\)，第8条缺项，即使 \(U=5\) 也不能补足。若仅将刑档改成三年至五年，则 \(L=3\) 进入边界，且其余四项成立，得「可以适用」。两案均须另外检查是否有独立管辖根据。另一个边界是运输工具仅过境。若境内既无犯罪行为亦无相关结果，且无船舶、航空器等独立依据，不能仅以运输路线推得 \(A\) 或 \(R\)。

\paragraph{时间效力\quad 从旧兼从轻}
令 \(g_o,g_n\) 表示依行为时旧法、后生效新法对同一行为是否构成犯罪，先固定行为时点和待比较版本。对尚未生效裁判的典型案件，在追诉等条件具备的前提下，第12条的决策可写为
\[
 \mathrm{Law}(x)=
 \begin{cases}
   \text{旧法},& \neg g_o,\\
   \text{新法},& g_o\land\neg g_n,\\
   \text{新法},& g_o\land g_n\land I_n\prec I_o,\\
   \text{旧法},& g_o\land g_n\land \neg(I_n\prec I_o).
 \end{cases}
\]
\(I_n\prec I_o\) 表示新法的适用刑档更轻。依最高人民法院关于适用刑法第12条的解释，先比较具体行为所适用档的法定最高刑，最高刑相同再比最低刑。构造同种主刑的例子，旧法为二年至五年，新法为三年至四年，则新法上限四年低于旧法上限五年，故适用新法。将新法刑档改为三年至五年，则两上限相同，新法下限更重，故适用旧法。两次比较均选择一套规则，不拼出二年至四年的第三套刑档。第12条还保留依当时法律已经作出的生效判决，不能直接把待决案件流程套到已生效裁判。

\subsection*{闭卷复述与检验}
\textbf{口诀}\quad「词分条，理由多。成事严定，地人保普。旧无止，双有比最高，平再比最低。」其中「成事严定」对应成文、事前、严格、确定，「地人保普」对应第6至9条。「旧无止」指旧法不认为是犯罪时对该行为不以新法入罪。

\textbf{复述路线}\quad
先写条文号和需解释的词，圈出通常外延与可能边界。再写入罪依据及各必要要素。跨境题画行为地、结果地、国籍、被害目标四格，境外案件再比上下限和当地可罚性。时间题画行为、变法、裁判三个时点，用旧法门槛排除事后增罪，再比较新旧适用刑档。

\begin{enumerate}
\item 只知暴力方式严重，能否把任意不同性质的对象纳入某罪名？\\
\textbf{答}\quad 不能。先验该对象是否进入可能文义，当然解释另须证明同类可比较及规则单调。缺少这些前提，「更严重」无法代入。
\item 普通中国公民在境外犯本法规定之罪，法定最高刑恰为三年有期徒刑。外国人另在境外对中国公民犯本法规定之罪，法定最低刑也恰为三年，且当地可罚。分别得什么？\\
\textbf{答}\quad 第一案由第7条适用本法，并进入可以不予追究的分支。第二案第8条各项满足，得到可以适用本法。两个「三年」比较的端点和法律效果均不同。
\item 旧法与新法都认为有罪，适用档分别是三年至七年、二年至七年，怎样选？\\
\textbf{答}\quad 最高刑同为七年，再比较最低刑。新法二年小于旧法三年，按较轻的新法。若旧法原本不入罪，则直接依旧法，不进入刑档比较。
\end{enumerate}
